% QSC v4.8.0-r01 candidate insertion
% Planck-Readout Metric Layer and MM/HC Field-Equation Refactor
% Internal research branch only. Not a public v4.7.0 mutation.

\section{Planck-Readout Metric Layer}
\label{sec:v480-r01-planck-readout-metric-layer}

\begin{quote}\small
\noindent\textbf{Status.} Internal v4.8.0-r01 research branch. This section is a
conservative readout-calibration layer. It does not alter the public v4.7.0
release, does not assert global T2 closure, and does not promote negative length
as a physical observable.
\end{quote}

\begin{definition}[Readout-calibrated effective Planck scale]
Let $R$ be an admissible readout/interface mode and let $\lambda_R$ denote the
local calibration factor induced by the MM/HC/QAM transport from the structural
interface to the physical readout. The effective Planck scale seen by $R$ is
\[
    \ell_{\mathrm{eff}}(R) := \ell_P\,\lambda_R .
\]
If $\lambda_R>0$, the readout preserves length orientation. If $\lambda_R<0$,
the sign is read as orientation, phase, or order reversal at the interface level,
not as a negative physical distance.
\end{definition}

\begin{remark}[Why the raw Planck ruler need not be invariant below physics]
The usual Planck length is a stabilized physical-readout scale. In the present
structure-first language, the lower information/interface layer need not carry
the same raw ruler. It carries a transport law whose physical projection may
reconstruct $\ell_P$ after calibration. Thus the invariant object is not the
unqualified number $\ell_P$ in every layer, but the guarded recovery of the
physical Planck scale under an admissible readout.
\end{remark}

\begin{lemma}[Matter readout recovery]
Assume that $R_{\mathrm m}$ is a stable matter readout and that its calibration
factor satisfies
\[
    0 < \lambda_{R_{\mathrm m}} = 1 + \varepsilon_{\mathrm m},
    \qquad |\varepsilon_{\mathrm m}| \leq \delta_{\mathrm m} .
\]
Then
\[
    \ell_{\mathrm{eff}}(R_{\mathrm m})
    = \ell_P(1+\varepsilon_{\mathrm m})
\]
recovers the ordinary Planck scale up to the declared local calibration error.
\end{lemma}

\begin{proof}
This is immediate from the definition of $\ell_{\mathrm{eff}}(R)$ and the
bounded calibration assumption. No global statement is used.
\end{proof}

\begin{lemma}[Opposite-orientation readout]
Assume that $R_{\tau}$ is a tachyon/opposite-orientation readout with
\[
    \lambda_{R_{\tau}} = -\mu_{R_{\tau}},
    \qquad \mu_{R_{\tau}}>0 .
\]
Then
\[
    \ell_{\mathrm{eff}}(R_{\tau}) = -\ell_P\mu_{R_{\tau}}
\]
is an oriented interface length. Its observable magnitude is
\[
    |\ell_{\mathrm{eff}}(R_{\tau})|=\ell_P\mu_{R_{\tau}},
\]
so the negative sign records orientation/phase/order reversal rather than a
negative physical distance.
\end{lemma}

\begin{proof}
The statement separates the signed interface coordinate from the positive metric
magnitude. The sign is therefore part of the readout grammar, while the physical
length magnitude remains non-negative.
\end{proof}

\begin{remark}[Connection to two-sided structural tensor notation]
Using the two-sided notation, the same calibration can be recorded schematically as
\[
    {}^{R}{}_{\mathrm{phys}}\mathcal T^{\mathrm{struct}}{}_{\mathrm{HC/MM}}
\]
where the left indices encode the readout-facing interface and the right indices
encode transport on the structure and action side. Ordinary upper--lower contractions
remain classical readouts; same-variance comparisons remain guarded by explicit
comparison tensors.
\end{remark}

\begin{remark}[No global closure claim]
This section does not prove exact through-collapse transport, exact all-readout
terminality, C6/T2 closure, or a global photon-clock/RH theorem. It only defines
a local readout-calibrated length grammar suitable for later MM/HC field-equation
refactoring.
\end{remark}
